\section{The Final Certificate}
\label{sec:certificate}

We apply Proposition~\ref{fw:transfer} to the fields of
Theorem~\ref{tw:field-family}, with the ceiling of
Theorem~\ref{an:ceiling} and the rational data below. Every parameter
written as a finite decimal is exact. The inequalities between decimals
displayed in this section, and the entries of
Tables~\ref{cert:tab-local} and~\ref{cert:tab-budget}, are computed and
printed by the supplementary program \texttt{geom241.py}
(Section~\ref{cert:computations}); the decimals are rounded outward from
interval enclosures and are never substituted into later calculations.

\subsection{Parameters and Local Data}
Set
\[
 \delta=\frac{4273}{100000}=\deltastar,\qquad p=\frac2{1+\delta}
 =\frac{200000}{104273},\qquad C=\frac{4871285}{10^8}=\ceiling,
\]
\[
 \theta_*=\frac{65535}{131072}.
\]
At the compact coordinates we use the Gaussian of
Lemma~\ref{pf:gaussian}, $f_\R^p=e^{-a_\R|z|^2}$ with
$a_\R=2p\delta=17092/104273<1$. At the pair coordinates we use the
profile $g_\C$ of Section~\ref{pf:section}, with the coefficients
\eqref{pf:bernstein-witness}, $s=1.14600585$ and
$a=0.0000128633241758$. For the Fourier bound we take
$\rho=10^{-3}$. For the mass we partition $[0,1]^2$ into the $64$
squares of side $1/8$ with rational corners and use the binomial degree
$N=18$ in Proposition~\ref{pf:mass}. For the overlap we use $m=1$ in
Proposition~\ref{pf:overlap}.

The selected rational primes are $\mathcal R=\{2,3,5,7,29\}$, with the
uniform local types of Theorem~\ref{tw:field-family}. At every place of
$F$ above $r\in\mathcal R$ we use the shell profile with product
weights $w_{ij}=w_{r,i}w_{r,j}$, in the notation of
\eqref{fw:product-shells} with $x=x'=(w_{r,i})_i$, and with
$Q=Q_r=r^{f_r}$ and $k=k_r$. Table~\ref{cert:tab-local} lists the local
data and Table~\ref{cert:tab-weights} the exact weights. At $7$ and $29$
the weights of the outer shells are zero, so these profiles have only
one and two nontrivial shells, and their gain over the hard window is
small (Table~\ref{cert:tab-local}). By \eqref{fw:product-shells}, the
local functional is
\begin{equation}\label{cert:finite-factor}
 \log\mathcal F_{\delta,r}
 =-\delta k_r\log Q_r+\log\bigl((k_r+1)S_r^2+2S_rR_r\bigr)
 -2(1+\delta)\log P_r,
\end{equation}
where $P_r=\sum_im_iw_{r,i}^p$, $S_r=\sum_im_iw_{r,i}^2$ and
$R_r=\sum_{i,i'}\Delta_{ii'}w_{r,i}w_{r,i'}$, with the shell masses
$m_i$ and the matrix $\Delta$ of Section~\ref{fw:section} for $Q=Q_r$.

\begin{table}[htbp]
\centering\small
\caption{Local data at the selected primes. Here $Q_r=r^{f_r}$, the
profile at $r$ uses the weights $w_{r,0},\dots,w_{r,n_r-1}$, and the
last two columns give $\log\mathcal F_{\delta,r}/(e_rf_r)$ for the shell
profile and for the hard window, rounded to ten decimals.}
\label{cert:tab-local}
\begin{tabular}{rrrrrrrr}
\toprule
$r$ & $e_r$ & $f_r$ & $Q_r$ & $k_r$ & $n_r$ & shells & hard window\\
\midrule
$2$ & $8$ & $4$ & $16$ & $7$ & $6$ & $0.0393401103$ & $0.0390666415$\\
$3$ & $2$ & $2$ & $9$ & $9$ & $6$ & $0.3675484302$ & $0.3643996093$\\
$5$ & $2$ & $2$ & $25$ & $6$ & $6$ & $0.2817099488$ & $0.2801636913$\\
$7$ & $1$ & $8$ & $5764801$ & $1$ & $2$ & $0.0034946608$ & $0.0034946569$\\
$29$ & $1$ & $4$ & $707281$ & $1$ & $3$ & $0.0294023321$ & $0.0294022443$\\
\bottomrule
\end{tabular}
\end{table}

\begin{table}[htbp]
\centering\small
\caption{The exact shell weights $w_{r,i}$ for $i\ge1$; all $w_{r,0}=1$.
The last column gives the weights rounded to six significant digits,
for orientation only; the certificate uses the exact fractions.}
\label{cert:tab-weights}
\begin{tabular}{rrlr}
\toprule
$r$ & $i$ & $w_{r,i}$ & $\approx w_{r,i}$\\
\midrule
2 & 1 & $2854206002066/59041214756005$ & $4.83426\cdot10^{-2}$\\
2 & 2 & $198349814036/95709219081499$ & $2.07242\cdot10^{-3}$\\
2 & 3 & $7762557303/86561984122975$ & $8.96763\cdot10^{-5}$\\
2 & 4 & $109894931/27828186078256$ & $3.94905\cdot10^{-6}$\\
2 & 5 & $17033757/96603874664399$ & $1.76326\cdot10^{-7}$\\
\midrule
3 & 1 & $5693537714695/63764117760919$ & $8.92906\cdot10^{-2}$\\
3 & 2 & $567838527095/78237129879259$ & $7.25792\cdot10^{-3}$\\
3 & 3 & $43356430379/72943614524859$ & $5.94383\cdot10^{-4}$\\
3 & 4 & $2265261068/45953937834745$ & $4.92942\cdot10^{-5}$\\
3 & 5 & $74149907/18018553422030$ & $4.11520\cdot10^{-6}$\\
\midrule
5 & 1 & $1343081168627/46386563301209$ & $2.89541\cdot10^{-2}$\\
5 & 2 & $61788974773/82067065665364$ & $7.52908\cdot10^{-4}$\\
5 & 3 & $1389672153/69896165011649$ & $1.98820\cdot10^{-5}$\\
5 & 4 & $32644248/60812698834313$ & $5.36800\cdot10^{-7}$\\
5 & 5 & $511565/34657937034136$ & $1.47604\cdot10^{-8}$\\
\midrule
7 & 1 & $2309805/69240611498956$ & $3.33591\cdot10^{-8}$\\
\midrule
29 & 1 & $36498535/95956876917664$ & $3.80364\cdot10^{-7}$\\
29 & 2 & $4/29090285629613$ & $1.37503\cdot10^{-13}$\\
\bottomrule
\end{tabular}
\end{table}


The weights are positive and nonincreasing, and the coefficients
$\beta_{ij}$ are positive; these are exact rational comparisons. The
rational number $Z(g_r)Q_r^{-k_r}$ is computed twice for each $r$: from
\eqref{fw:product-shells}, and by writing the profile as the
nonnegative combination
$\sum_{i,j}(w_{r,i}-w_{r,i+1})(w_{r,j}-w_{r,j+1})\mathbf 1_{\mathcal
B_i\times\varpi^{-k_r}\mathcal B_j}$ of indicators of products of balls,
with $w_{r,n_r}=0$, and applying \eqref{fw:rectangle-gram}. The two
exact values agree.

\subsection{Certified Enclosures}
Let $q=sp-1=12492817/10427300$, and let $T$ and $T'$ be independent
with density $qt^{q-1}$ on $(0,1)$. Proposition~\ref{pf:mass}, with the
data above, gives
\begin{equation}\label{cert:mass}
 38.7888697256861058<\bar A_\C=\mathbb E\,P(T,T')^p<38.7888697256861071 .
\end{equation}
Here every local Bernstein coefficient is positive, every local ratio
$\rho_{\mathcal Q}$ is less than $0.332$, and the tail allowances in
\eqref{pf:subrectangle-mass} add up to less than $5.9\cdot10^{-16}$. The
threshold conditions \eqref{pf:cross-threshold} hold for every cross,
and Proposition~\ref{pf:overlap} with $m=1$ gives
\begin{equation}\label{cert:overlap}
 348.4186894704059951<Z_*<348.4186894704059952 ,
\end{equation}
where $Z_*$ is the right side of \eqref{pf:collected-overlap}; its
signed remainder $E_{-,1}$ is less than $6.1\cdot10^{-17}$. The upper
endpoint in \eqref{cert:overlap} bounds this expression, not the
overlap $I_\C$. Let $J_\C^*$ be the right side of
\eqref{pf:collected-functional} with $\bar A_\C^*$ the upper endpoint
of the enclosure \eqref{cert:mass}. By Corollary~\ref{pf:functional},
$J_\C\ge J_\C^*$, and
\[
 1.3556525171042125<J_\C^*<1.3556525171042127 .
\]
Formula~\eqref{pf:gaussian-compact} gives
\[
 -0.2107595968649958<J_\R<-0.2107595968649957 .
\]
For the finite places, \eqref{cert:finite-factor} and the data of
Table~\ref{cert:tab-local} give
\[
 0.7214954821109038790<\sum_{r\in\mathcal R}
 \frac{\log\mathcal F_{\delta,r}}{e_rf_r}<0.7214954821109038791 ,
\]
against $J-\delta H<0.7165268434$ for the hard windows. We also have
\begin{gather*}
 15.6917787983405342<H<15.6917787983405343,\\
 5.6560047235563094<\ell<5.6560047235563095 ,
\end{gather*}
where $H=\sum_rk_r\log r/e_r$ and
$\ell=\frac94\log2+\frac12\log3615$.

For the Fourier condition, the rational tube quantity of
Lemma~\ref{pf:tube} satisfies $q_0<0.026<1$, and
\[
 \log K_\C<0.0491254364,\qquad \sigma_\C>1.7518755886,\qquad
 \mu=2e^{H/2-\ell}>17.8683654522 .
\]
Since $K_\C<4$ and $\sigma_\C>1$, Proposition~\ref{pf:admissible} shows
that $(f_\R,g_\C)$ is admissible with Fourier constants $M=2$ and
$\sigma=1$. The right side of \eqref{geo:poisson-condition} is then
$\log50+2<5.9121$, so \eqref{fw:fourier-condition}, which by
\eqref{fw:uniform} is the same condition with $H_f=H$, holds with a
wide margin: although $\mu$ is not large, only
$\mu\ge\log M+2\log5+2$ is needed.

Define the conservative margin
\[
 \begin{split}
 M_*(\theta)={}&\sum_{r\in\mathcal R}
 \frac{\log\mathcal F_{\delta,r}}{e_rf_r}
 -(1/2-\delta)\ell-C+(1-\delta)\log2+(1-\theta)\log\pi\\
 &+(1-2\theta)J_\R+\theta J_\C^* .
 \end{split}
\]
By \eqref{fw:uniform}, $J_\C\ge J_\C^*$ and $\theta\ge0$, it is a lower
bound for \eqref{fw:margin}. It is affine in $\theta$ with slope
\[
 0.6324418249848039<-\log\pi-2J_\R+J_\C^*<0.6324418249848040 ,
\]
so on $[\theta_*,1/2)$ its minimum is at $\theta_*$, and
\begin{equation}\label{cert:margin}
 M_*(\theta_*)>0.000176730033534598 .
\end{equation}
Table~\ref{cert:tab-budget} lists the terms of $M_*(\theta_*)$. The
proof of Proposition~\ref{fw:transfer} only needs some $\eta>0$ with
$4\eta$ below the margin. For illustration, $\eta=10^{-9}$ gives
\begin{equation}\label{cert:final-margin}
 M_*(\theta_*)-4\eta>0.000176726033534598>0 .
\end{equation}
The margin is small compared with the individual terms. Since
$M_*$ decreases with slope $1$ in $C$ and with slope $1/2-\delta$ in
$\ell$, it would be absorbed by an increase of $1.77\cdot10^{-4}$ in $C$,
or of $3.87\cdot10^{-4}$ in $\ell$. With these profiles and shell
weights optimized in the same way for $\delta=0.0428$, the same
evaluation gives $M_*(\theta_*)<-0.0015$.

\begin{table}[htbp]
\centering\small
\caption{The terms of $M_*(\theta_*)$, rounded to ten decimals.}
\label{cert:tab-budget}
\begin{tabular}{lr}
\toprule
Term & Value\\
\midrule
$\sum_r\log\mathcal F_{\delta,r}/(e_rf_r)$ & $0.7214954821$\\
$-(1/2-\delta)\ell$ & $-2.5863212799$\\
$-C$ & $-0.0487128500$\\
$(1-\delta)\log2$ & $0.6635290015$\\
$(1-\theta_*)\log\pi$ & $0.5723736765$\\
$(1-2\theta_*)J_\R$ & $-0.0000032159$\\
$\theta_*J_\C^*$ & $0.6778159157$\\
\midrule
$M_*(\theta_*)$ & $0.0001767300$\\
\bottomrule
\end{tabular}
\end{table}

\begin{remark}\label{cert:signature-remark}
The certificate uses the bound $\theta\ge\theta_*$ only through the
positivity of $M_*$ on $[\theta_*,1/2)$. The program
\texttt{geom241.py} also encloses the value at $\theta=0$ and the zero
$\theta_0$ of the affine function $M_*$:
\[
 M_*(0)<-0.316,\qquad \theta_0<0.49971293 .
\]
So $M_*(\theta)>0$ for $\theta\ge0.49971293$. By the proof of
Theorem~\ref{tw:field-family}, $b/d=1/(2N_\iota)$, where $N_\iota$ is
the number of conjugates of $\iota_1$ in $\Gal(K/B)$, so that
$\theta=1/2-1/(4N_\iota)$. Hence $N_\iota\ge871$ would suffice, whereas
that proof gives $N_\iota\ge2^{15}$. On the other hand, at $\theta=0$ the margin
\eqref{fw:margin} does not involve $J_\C$ and equals $M_*(0)<0$: with
these profiles, mixed signature is essential.
\end{remark}

\subsection{Proof of the Main Theorem}

\begin{proof}[Proof of Theorem~\ref{thm:main}]
For each $n\ge49$, choose a Galois extension $K_n/B$ of degree $2^n$ as
in Theorem~\ref{tw:field-family} (with $m=2^n$), a complex conjugation
$\iota_1$ at a place of $K_n$ above $v_1$, and put
$F_n=K_n^{\langle\iota_1\rangle}$. Then $d_n=[F_n:\Q]=2^n\to\infty$. By
that theorem, each $K_n/F_n$ satisfies (G1)--(G3) with
$\lambda=2^{9/4}\sqrt{3615}$, has the uniform local types
$(e_r,f_r)_{r\in\mathcal R}$ of Table~\ref{cert:tab-local}, and has
$\theta_*\le\theta_n<1/2$. By Theorem~\ref{an:ceiling},
$d_n^{-1}\log L_{F_n}(1)<C$ for all sufficiently large $n$; we discard
the finitely many other fields.

Let $\mathcal P_n$ be the set of places of $F_n$ above $\mathcal R$,
with the shell profiles of Table~\ref{cert:tab-weights}. Their types
form a fixed set of five. By Proposition~\ref{pf:admissible} and the
enclosures above, $(f_\R,g_\C)$ is admissible with Fourier constants
$M=2$, $\sigma=1$, and \eqref{fw:fourier-condition} holds for every
$n$. By \eqref{fw:uniform}, Corollary~\ref{pf:functional} and the slope
bound, the margin \eqref{fw:margin} of $K_n/F_n$ is at least
$M_*(\theta_n)\ge M_*(\theta_*)$, which is positive by
\eqref{cert:margin}. Proposition~\ref{fw:transfer} therefore gives finite
sets $U_n\subset\R^2$ with $|U_n|\to\infty$ and
$u(U_n)/|U_n|^{1+\delta}\to\infty$, where $1+\delta=1.04273$; these are
the sets of Theorem~\ref{thm:main}, indexed by $n$. In particular
$u(U_n)\ge|U_n|^{1.04273}$ for all large $n$, which gives the last
assertion of Theorem~\ref{thm:main} at the cardinalities $|U_n|$.
\end{proof}

\subsection{Computations}\label{cert:computations}
The program \texttt{geom241.py} evaluates the formulas of
Sections~\ref{fw:section} and~\ref{pf:section} for the data above: the
local functionals \eqref{cert:finite-factor} from the exponents and
weights in the data file \texttt{shells241\_0.04273.json}, and the
quantities of Propositions~\ref{pf:mass} and~\ref{pf:overlap}, of
Lemma~\ref{pf:tube} and of Corollary~\ref{pf:functional} from the
profile data above. For these it uses, unchanged, three certified
modules of an earlier supplementary archive of the author
(Section~\ref{intro:code}), which implement exactly these formulas. It checks
every finite hypothesis used above: the positivity and monotonicity of
the weights, the two exact computations of each $Z(g_r)Q_r^{-k_r}$, the
positivity of the $\beta_{ij}$ and of the local Bernstein coefficients,
the conditions \eqref{pf:cross-threshold}, $q_0<1$, $a_\R\le1$,
$\sigma_\C>1$, $\log K_\C<2\log2$, $\mu>10$ and the positivity of the
slope. It then prints the enclosures stated in this section. The
exponents and weights were produced by the separate optimization
program \texttt{shells241.py}; they enter the certificate only as the
exact rational data of Table~\ref{cert:tab-weights}.

Rational quantities are computed exactly: $q_0$, the coefficients of
the powers $W_{\mathcal Q}^k$ and the binomial coefficients in
Proposition~\ref{pf:mass}, $G_n$ and $\widetilde G_n$, and the shell sums
$S_r$ and $R_r$. The remaining quantities involve the logarithm, the
exponential
(also through the real powers $x^y=\exp(y\log x)$ in
\eqref{pf:subrectangle-moments}, in $c_{\mathcal Q}^p$ and in
$w_{r,i}^p$), the square root in $\sigma_\C$,
the constant $\pi$, and the digamma function, which enters only through
Lemma~\ref{pf:digamma}. They are evaluated in the interval arithmetic of
mpmath \cite{Mpmath}, which computes each arithmetic operation and each
of these elementary functions with the lower endpoint rounded down and
the upper endpoint rounded up, at a working precision of $80$ decimal
digits. Hence every computed interval contains the exact value, and the
decimals stated above, which are coarser than the working precision by
more than fifty orders of magnitude, are rigorous bounds.
